\begin{abstract}
Model ensembling and weight-space merging have emerged as powerful paradigms for
integrating multi-task experts without the prohibitive computational costs of
joint training. However, when adapting ensembling coefficients across layers on
small calibration datasets, the resulting high-dimensional search space is
highly susceptible to transductive overfitting. In this work, we approach this
challenge through the lens of statistical learning theory and introduce
\textbf{Rademacher-Bounded Fourier Trajectory Merging (RB-FTM)}. Instead of
treating layer-wise ensembling coefficients as independent parameters, we
project them onto a low-frequency harmonic Fourier subspace across network
depth. We derive a novel, mathematically rigorous empirical Rademacher
complexity bound for this trigonometric ensembling class, proving that its
trajectory-space capacity (fluctuation complexity across layers) is strictly
bounded by the spectral cutoff frequency $F$ and network depth $L$, independent
of the parameter count of the underlying neural network. To physically enforce
this complexity bound, we introduce an analytical Spectral Lasso ($L_1$)
regularizer on the Fourier coefficients. Furthermore, we show that our
trigonometric formulation naturally mitigates the severe boundary runaway
characteristic of prior polynomial ensembling trajectories at early feature-
extraction and final classification layers. Empirical evaluations inside a
mathematically controlled, synthetic \textbf{Analytical Coordinate Sandbox
(ACS)}—which simulates representation propagation through the layer depth
coordinates of Deep12LayerCNN and CLIP ViT-B/16 backbones—demonstrate that while
the parameter-free Static Uniform baseline remains a highly robust consensus
upper bound in perfectly coordinate-aligned spaces, RB-FTM and its Discrete
Cosine Transform (DCT) variant, RB-DCTM, significantly outperform all other
adaptive and tuned ensembling baselines, achieving up to $70.70\%$ accuracy on
CNN and $72.70\%$ on CLIP, while generating smooth, stable, and theoretically
grounded parameter trajectories.
\end{abstract}
